% Generated draft: requires mathematical and source review.
\documentclass[11pt]{article}
\usepackage[margin=25mm]{geometry}
\usepackage{fontspec}
\setmainfont{Latin Modern Roman}
\usepackage{amsmath,amssymb,mathtools,unicode-math}
\setmathfont{Latin Modern Math}
\usepackage{enumitem,needspace,xurl,hyperref}
\setlength{\parindent}{0pt}
\setlength{\parskip}{5pt}
\setlength{\emergencystretch}{3em}
\newcommand{\N}{\mathbb{N}}
\newcommand{\Z}{\mathbb{Z}}
\newcommand{\Q}{\mathbb{Q}}
\newcommand{\R}{\mathbb{R}}
\newcommand{\C}{\mathbb{C}}
\newcommand{\F}{\mathbb{F}}
\newcommand{\A}{\mathbb{A}}
\newcommand{\PP}{\mathbb{P}}
\newcommand{\E}{\mathbb{E}}
\newcommand{\eps}{\varepsilon}
\newcommand{\dd}{\,\mathrm{d}}
\begin{document}
\section*{UnsolvedMath ID: 30004150}
\textbf{Source problem code:} catalogue entry
\subsection{Definitions and mathematical statement}
Let \(a\) denote the causal regression coefficient vector referred to in the supplied record, and let \(\hat a\) denote an estimated regression vector in the same coefficient space. Write
\[
\Delta:=\hat a-a
\]
for their discrepancy. The record identifies finite-sample noise and confounding as possible causes of this discrepancy. Here confounding refers to distortion of an observational association by common causes of the variables under study; the record specifies no causal model that makes this description operational. Regularization refers to a restriction or penalty used in estimating the regression vector; its form is not specified.

\textbf{Research question.} When \(\Delta\) arises from finite-sample noise or from confounding, determine what justification there is for algorithms to distinguish these causes, and whether regularization can mitigate confounding in a manner analogous to its mitigation of overfitting.

\textbf{Unresolved formulation.} The coefficient domain and dimension, regression model, sampling assumptions, confounding mechanism, causal identification assumptions, algorithm class, regularization scheme, and criterion for mitigation are not supplied. In particular, the record does not specify whether mitigation means improved recovery of \(a\), improved prediction under intervention, or another objective. No identifiable decomposition of \(\Delta\) into noise and confounding components is asserted. Consequently, this is a research question rather than a fully quantified conjecture; a precise theorem or optimization target requires additional hypotheses.

The supplied literature assessment describes results under an explicit multivariate confounder model, including a corresponding generalization bound, while describing the question beyond those assumptions as unresolved. Those model assumptions and the bound are not reproduced in the record and are not asserted here as independently verified facts.
\subsection{Short English statement}
The question asks why a learning algorithm should treat estimation noise and confounding differently when both can move an estimated regression vector away from its causal target, and whether regularization can reduce confounding effects. A precise guarantee requires a specified causal model and a specified measure of success. The supplied status is partially solved, subject to dated curation and human mathematical review.
\subsection{Sources}
\begin{itemize}
\item Local UnsolvedMath registry, permanent ID \(30004150\), title \emph{Distinguishing Confounding from Estimation Noise}, source problem code \texttt{catalogue entry}, in the Oberwolfach Reports Open Problems collection. No individual proposer is supplied.
\item Supplied report citation: \emph{Foundations and new horizons for causal inference} (2020), Oberwolfach Reports, pp.\ 1499--1571, DOI \texttt{10.4171/owr/2019/25}. \url{https://doi.org/10.4171/owr/2019/25}. The registry records proposed year 2019; both dates are retained as supplied.
\item Supplied literature source labeled ``Janzing'': \url{https://arxiv.org/abs/1906.12179}.
\item Supplied literature source labeled ``NeurIPS version'': \url{https://papers.nips.cc/paper/9432-causal-regularization}.
\item Literature assessment checked at \texttt{2026-08-22}: supplied dated curation, not independently verified. It labels the problem partially solved and attributes to causal regularization an explanation of the resemblance between confounding and overfitting in an explicit multivariate confounder model, together with a corresponding generalization bound. It states that the broad question outside those assumptions is not fully resolved and requests independent verification of time-sensitive status and cited claims.
\end{itemize}
\end{document}
